\chapter{Future Work}
    The following are some of the possibilities of continuation to my current work in the area of Static and Dynamic Gomory-Hu Tree and their Applications, for the upcoming Semester/Year for MTP (M.Tech Project):-

    \vspace{5mm}
    \begin{enumerate}
        \item {
            Further exploring Cut-based Image Segmentation and trying to improve for better execution times and overall more accurate final Clusters.
        }
        \item{
            Extending Dynamic Gomory-Hu Tree to multiple edge-updates instead of single-update, given that the updates are Offline i.e. known to us beforehand.
        }
        \item {
            Exploring and Improving Dynamic Gomory-Hu Tree for specific special graph types. Instead of going into improving Dynamic Gomory-Hu Tree for general graphs as a whole, I can take some special graph type and then try to improve Dynamic Gomory-Hu Tree for that specific case.
        }
        \item {
            Theoretically determining beforehand, the number of newer flow computations required in case of decreasing edge weights in Dynamic Gomory-Hu Tree. \\
            The current Decreasing Edge Weight algorithm for Dynamic Gomory-Hu Tree as proposed by T. Hartmann and D. Wagner (2013) [HT13] only gives an upper bound as to how many newer flow computations the algorithm will perform, the exact value is known only after fully performing the updates. \\
            Instead of this, I can attempt to mathematically/theoretically determine the number of newer flow computations that will be done by this algorithm without actually doing the update.
        }
        \item {
            Finally, I can explore other Applications of Gomory-Hu Tree apart from Image Segmentation such as - Network Reliability and Sensitivity Analysis or Community Detection.
        }
    \end{enumerate}